% ============================================================
% ТЕРМИНЫ И ОПРЕДЕЛЕНИЯ
% ============================================================
\chapter*{ТЕРМИНЫ И ОПРЕДЕЛЕНИЯ}
\addcontentsline{toc}{chapter}{Термины и определения}

В настоящей работе применяются следующие термины с соответствующими определениями.

\begin{sloppypar}

\textbf{Инференс} (англ. \textit{inference}) --- процесс получения предсказаний от обученной модели машинного обучения по новым входным данным. Отличается от процесса обучения: параметры модели зафиксированы, и вся вычислительная нагрузка сводится к прямому проходу через граф вычислений.

\textbf{Холодный старт} (англ. \textit{cold start}) --- состояние, при котором инстанс сервиса инференса создаётся с нуля: контейнер не был запущен, кэш весов отсутствует, и перед обработкой первого запроса необходима полная инициализация, включая загрузку файлов модели.

\textbf{LLM} (Large Language Model, большая языковая модель) --- языковая модель на базе архитектуры Transformer с числом параметров от нескольких миллиардов и выше, требующая для хранения весов от десятков до сотен гигабайт.

\textbf{Веса модели} --- набор числовых параметров нейросети (матриц, векторов смещений), определяющих поведение модели. Для современных LLM хранятся в форматах \texttt{.safetensors} или \texttt{.bin} и загружаются в GPU-память перед инференсом.

\textbf{Чанк} (англ. \textit{chunk}) --- фрагмент файла весов модели фиксированного размера, используемый как минимальная единица кэширования и P2P-передачи в разработанной системе.

\textbf{Serverless} (бессерверные вычисления) --- модель развёртывания, при которой платформа автоматически управляет выделением и освобождением вычислительных ресурсов в зависимости от нагрузки, вплоть до полного удаления инстансов при её отсутствии (scale-to-zero).

\textbf{Kubernetes} --- система оркестрации контейнеров с открытым исходным кодом, обеспечивающая автоматизацию развёртывания, масштабирования и управления контейнеризованными приложениями.

\textbf{Pod} (под) --- минимальная единица развёртывания в Kubernetes, объединяющая один или несколько контейнеров, разделяющих сетевое пространство имён и тома хранилища.

\textbf{DaemonSet} --- ресурс Kubernetes, гарантирующий запуск ровно одного экземпляра указанного пода на каждой worker-ноде кластера (или на подмножестве нод, определяемом селектором).

\textbf{Scale-to-zero} --- способность платформы снижать число запущенных реплик сервиса до нуля при отсутствии входящих запросов, обеспечивая экономию ресурсов.

\textbf{SLA} (Service Level Agreement) --- соглашение об уровне обслуживания, определяющее количественные показатели качества сервиса, в том числе максимально допустимые задержки обработки запросов.

\textbf{P2P} (peer-to-peer, одноранговая сеть) --- сетевая архитектура, при которой узлы напрямую обмениваются данными без центрального сервера. В контексте диплома применяется для передачи чанков между \texttt{storage-agent} на разных нодах.

\textbf{Control Plane} --- плоскость управления; в контексте работы --- Redis-кластер, хранящий метаданные о местоположении чанков и состоянии агентов.

\textbf{Data Plane} --- плоскость данных; в контексте работы --- совокупность локальных дисков worker-нод, по которым распределены чанки, и каналы их доставки (UDS, TCP, HTTP к внешнему репозиторию).

\textbf{Redis tracker / metadata registry} --- использование Redis в роли реестра метаданных: каждый \texttt{storage-agent} периодически публикует в Redis информацию о своём IP-адресе, доступных портах и хранящихся чанках с TTL-меткой.

\textbf{FUSE} (Filesystem in Userspace, файловая система в пространстве пользователя) --- механизм ядра Linux, позволяющий реализовать файловую систему в виде обычного пользовательского процесса. Системные вызовы VFS перенаправляются в пользовательский демон через ядерный модуль FUSE.

\textbf{UDS} (Unix Domain Socket, сокет домена Unix) --- механизм IPC (межпроцессного взаимодействия), обеспечивающий обмен данными между процессами в пределах одного узла через файловую систему без IP-стека.

\textbf{NFS} (Network File System) --- сетевой протокол распределённой файловой системы, позволяющий монтировать удалённые файловые системы через сеть и получать к ним доступ как к локальным.

\textbf{flock} --- системный вызов POSIX для установки рекомендательных (advisory) файловых блокировок уровня процесса; применяется в работе для координации параллельных загрузок модели в NFS-решении.

\textbf{Lease} --- ресурс Kubernetes API (\texttt{coordination.k8s.io/v1/Lease}), предоставляющий временной токен владения ресурсом; используется механизмом Leader Election для выбора единственного активного загрузчика.

\textbf{Leader Election} --- алгоритм распределённого консенсуса, гарантирующий, что в каждый момент времени среди конкурирующих процессов ровно один является лидером и выполняет привилегированное действие (например, скачивание модели).

\textbf{Cache hit} (попадание в кэш) --- ситуация, при которой запрошенные данные обнаружены в кэше и могут быть возвращены без обращения к более медленному источнику.

\textbf{Cache miss} (промах кэша) --- ситуация, при которой запрошенные данные отсутствуют в кэше и необходимо обращение к следующему уровню иерархии хранения.

\textbf{Fallback} --- механизм переключения на резервный источник данных при недоступности основного; в контексте диплома --- загрузка чанка из внешнего репозитория HuggingFace при отсутствии его в локальном и соседних кэшах.

\textbf{SPOF} (Single Point of Failure, единственная точка отказа) --- компонент системы, выход из строя которого вызывает полную недоступность всей системы или её критически важной части.

\end{sloppypar}
